\clearpage
\section{Anhang: Nutzung der Simulation}

Die \href{https://github.com/danielschwalbe/ipres}{Simulation} ist ein in der
\href{https://www.python.org}{Python}-Programmiersprache geschriebenes Projekt. Es enthält neben
der \textbf{ipres} Python-Package viele Readmes und Beispielnotebooks.
\href{https://jupyter.org}{Jupyter Notebooks} ist eine aus der Python-Welt stammende Technologie,
die es erlaubt, Markdown (Dokumentation) und Code in einer Datei zusammenzufassen und interaktiv
auszuführen. Das macht sie geeignet, um z.\,B.\ eben solche Wahlverfahren zu demonstrieren. Man
kann Parameter ändern und live sehen, wie sie sich auswirken.

Eigentlich müsste man einen Jupyter-Lab-Server starten, um die Notebooks auszuführen. Das Projekt
verlinkt aber zu \href{https://mybinder.org/v2/gh/danielschwalbe/ipres/HEAD}{binder}, welches das
Projekt in eine virtuelle Maschine kopiert und dort den Jupyter-Server bereits ausführt. Man muss
nur ein bisschen Geduld haben, bis die Maschine gestartet ist. Bis dahin sieht es so aus:

\begin{figure}[H]
  \centering
  \includegraphics[width=0.85\textwidth]{screenshots/binder_start.png}
  \caption{Binder lädt die virtuelle Maschine}
\end{figure}

Wenn die Maschine hochgefahren ist, wird man von folgendem Bildschirm begrüßt:

\begin{figure}[H]
  \centering
  \includegraphics[width=0.85\textwidth]{screenshots/binder_loaded.png}
  \caption{JupyterLab nach dem Start}
\end{figure}

Um zu den Notebooks zu gelangen, muss man auf \texttt{notebooks} klicken.

\begin{figure}[H]
  \centering
  \includegraphics[width=0.85\textwidth]{screenshots/binder_loaded_language.png}
   \caption{Notebooks-Ordner}
\end{figure}

\clearpage

Und dann auf \texttt{de}, um zu den deutschen Notebooks zu gelangen.

\begin{figure}[H]
  \centering
  \clipbox{0pt {\height-7cm} 0pt 0pt} {\includegraphics[width=0.85\textwidth]{screenshots/notebooks_overview.png}}
  \caption{Deutsche Notebooks}
\end{figure}

Hier kann man sich dann welche herauspicken, z.\,B.\ die \texttt{bundestagswahl.ipynb}.

Durch Klick auf den Startpfeil wird die gerade markierte Zelle ausgeführt. Man beachte, dass die
Zellen in der Regel von oben nach unten ausgeführt werden müssen, weil die unteren Zellen
meistens (nicht immer) die Ergebnisse der oberen Zellen voraussetzen.

\begin{figure}[H]
  \centering
   \clipbox{0pt {\height-7cm} 0pt 0pt}{\includegraphics[width=0.85\textwidth]{screenshots/example_run_cell.png}}
  \caption{Ausführung einer Notebook-Zelle}
\end{figure}

\clearpage
Mithilfe der zwei Pfeile kann man auch das ganze Notebook auf einmal ausführen:

\begin{figure}[H]
  \centering
  \clipbox{0pt {\height-7cm} 0pt 0pt}{\includegraphics[width=0.85\textwidth]{screenshots/example_run_all.png}}
  \caption{Alle Zellen auf einmal ausführen}
\end{figure}

Dann wird man gefragt, ob man den Kernel neu starten möchte.

\begin{figure}[H]
  \centering
  \includegraphics[width=0.85\textwidth]{screenshots/restart_kernel.png}
  \caption{Kernel neu starten}
\end{figure}

Solange man am Notebook nichts geändert hat, kann man hier immer ja sagen. Nur wenn man durch fortgesetztes Experimentieren
die Variablen in einen Zustand bekommen hat, der sich aus dem Notebookcode nicht rekonstruieren lässt, muss man wissen, dass
das Ergebnis nach dem Kernelneustart weg sein wird.

Wenn man am Code etwas ändern will, kann man die Zeilen einfach editieren und nochmal ausführen.
Man könnte z.\,B.\ in der oberen Zelle andere Seeds eintragen und würde sehen, wie andere Parteien
gewinnen.

Falls man weitergehende Änderungen machen möchte, muss man deswegen heutzutage nicht unbedingt
Python können. Die KIs können das auch und helfen gerne.
